% Reviewer-response working copy; figures/style/bibliography require the original project.
% !TEX program = pdflatex
% CVPR 2026 Introduction, Related Work, and Methodology draft; official style files are unmodified.
\documentclass[10pt,twocolumn,letterpaper]{article}
\usepackage[review]{cvpr}
\usepackage{graphicx}
\usepackage{amsmath}
\usepackage{amssymb}
\usepackage{booktabs}
\usepackage{cuted}
\usepackage{caption}
\definecolor{cvprblue}{rgb}{0.21,0.49,0.74}
\usepackage[pagebackref,breaklinks,colorlinks,allcolors=cvprblue]{hyperref}
\def\paperID{*****}
\def\confName{CVPR}
\def\confYear{2026}
\clubpenalty=10000
\widowpenalty=10000
\title{Sparse Specialist Feedback for Spatial Test-Time Adaptation
	in Spatio-Temporal Video Grounding}
\author{Anonymous authors}
\newcommand{\mvIoU}{\ensuremath{m_{\mathrm{vIoU}}}}
\newcommand{\mtIoU}{\ensuremath{m_{\mathrm{tIoU}}}}
\newcommand{\msIoU}{\ensuremath{m_{\mathrm{sIoU}}}}

\begin{document}
\AddToHookNext{cmd/@maketitle/after}{%
  \par\noindent
  \begin{minipage}{\textwidth}
  \captionsetup{type=figure}
  \centering
  \includegraphics[width=\textwidth]{figures/motivation.pdf}
  \caption{\textbf{Temporal--spatial mismatch in STVG.}
  For the same query, the three models obtain high temporal IoU (tIoU) but low spatial IoU (sIoU), illustrating that event localization and referent localization can diverge.
  Green dashed boxes and green timelines indicate ground truth; coral boxes and timelines indicate predictions.
  The scores on the right correspond to this example.}
  \label{fig:motivation}
  \end{minipage}\par
}
\maketitle

\begin{abstract}
Spatio-temporal video grounding must recover an event interval and the trajectory
of a language-specified instance under deployment shifts. We study sparse
specialist-guided spatial test-time adaptation through a 1,792-parameter
decoder interface. On-policy localization actions are scored against one
frozen grounding specialist, and their reward-weighted distribution guides
the current-query spatial update while the native temporal interval is retained.
The fixed system increases mean video IoU from 20.017 to 22.583 on HC-STVG-v2
and from 11.129 to 14.097 on VidSTG. Common-state comparisons distinguish the
system gain from the spatial objective: they do not establish stable OPD
superiority over direct regression, and output-budget controls reveal
dataset-dependent tradeoffs. A cached temporal--spatial factorial study finds
that an existing actionness projection adds 3.823 and 2.079 percentage points
to OPD on two historical 32-parent panels, with no established OPD-specific
interaction. These results identify a reproducible compact adaptation system
and clarify which benefits are attributable to sparse feedback, spatial
readout and temporal coverage.
\end{abstract}

\section{Introduction}
\label{sec:introduction}

Spatio-temporal video grounding (STVG) aims to localize the instance
referred to by a natural-language query in a video, returning
the interval in which the described event occurs and a sequence of
bounding boxes tracing its participant~\cite{vidstg,tubedetr}.
The task links language understanding to event localization and object
tracking, with applications in video retrieval and embodied
perception~\cite{tastvg}. Recent advances have improved joint
space--time decoding~\cite{tubedetr}, target-aware query
representations~\cite{tastvg}, and collaborative reasoning over
appearance and motion~\cite{dualstream,e3m}. Nevertheless, a model
trained on one dataset may encounter markedly different scenes,
appearances, motion patterns, and linguistic expressions after
deployment. Video corruptions, including missing or frozen frames,
blur, occlusion, and exposure changes, further alter the evidence on
which grounding depends~\cite{tadcorruptions,vitta}.

Adapting such a model through additional supervised training requires
newly annotated target videos, including event boundaries and
framewise boxes. These annotations are costly to obtain whenever the
deployment distribution changes. Test-time adaptation (TTA) instead
updates a pretrained model using incoming unlabeled
inputs~\cite{tent,cotta}, allowing its predictions to respond to
target-domain observations without a new annotation stage. We study
this setting for STVG, focusing on cross-dataset transfer and video
corruptions, and allow a frozen expert
to supply sparse spatial observations during inference.


Test-time adaptation relies on measurable surrogate objectives to update
a model without target annotations. In classification-oriented settings,
such as image classification, time-series recognition, and video action
recognition, adaptation typically operates on a class distribution or
its underlying representation. Entropy minimization, prediction
consistency, and feature alignment consequently provide practical
self-supervised signals~\cite{tent,cotta,vitta}. Although these signals
are imperfect, they offer a common basis for optimizing model
predictions under distribution shifts~\cite{pitfalls}. Grounding tasks,
however, demand more than confidence in a predicted category.
In image-based visual grounding (VG), the model must establish a
correspondence between a referring expression and a particular visual
region~\cite{mdetr,groundingdino}. Detecting a semantically relevant
object does not necessarily identify the intended referent, especially
when multiple instances share similar appearances or attributes.
Consequently, confidence in object detection and consistency of
predicted regions provide only partial evidence of grounding
correctness.

The challenge becomes more involved when grounding extends from
images to videos. Spatial video grounding (SVG) requires identifying
and maintaining the referred instance across frames, where changes
in appearance, occlusion, and visually similar distractors complicate
referent association~\cite{tempcd}. Temporal video grounding (TVG),
by contrast, seeks the interval in which a language-described event
occurs, requiring temporal relevance and event boundaries to be
distinguished from surrounding video content~\cite{videoopd}.
Spatio-temporal video grounding (STVG) unifies these requirements:
the model must identify the correct instance, determine when the
described event occurs, and localize that instance throughout the
corresponding interval~\cite{vidstg,tubedetr,e3m}. These decisions
are interdependent, yet the evidence supporting them is heterogeneous.
A correctly localized object may remain visible before or after the
queried event, while an accurately identified temporal interval may
still contain several plausible target instances. Hence, correctness
of an individual spatial or temporal cue does not establish
correctness of the complete spatio-temporal tube.
Figure~\ref{fig:motivation} illustrates this mismatch: accurate temporal
localization can coexist with spatial predictions that miss the referred target.
Under test-time
distribution shifts, this coupling makes adaptation particularly
challenging: without target annotations, the model must learn from
partial and potentially unreliable evidence while distinguishing
localization, temporal relevance, and referent identity. A central
challenge for STVG-TTA is therefore how to convert such heterogeneous
test-time evidence into effective corrections of a jointly structured
grounding prediction.


Pretrained image grounding models offer a practical source of
frame-level localization evidence~\cite{groundingdino}. Their
predictions can help localize an instance, while the surrounding
event and its temporal boundaries remain properties of the video
query. We therefore use the STVG model's native temporal interval to
select where spatial observations are collected and retain that
interval as the temporal component of its output. Within the selected
event, a frozen Grounding DINO model supplies reference boxes on a
small number of frames. These references provide localized feedback
for adapting the student's spatial predictions over the video.

A frame-level specialist reference specifies one localization
hypothesis, but the student's prediction can be refined in several
nearby directions. Direct box regression treats the reference
coordinates as a fixed target; it does not explicitly compare these
alternatives as the student's output changes. Moreover, pseudo-boxes
from pretrained vision--language detectors can contain localization
noise, as documented in open-vocabulary detection~\cite{sasdet}.
On-policy distillation provides another learning interface: it
obtains feedback on the current student's generated outputs and uses
that feedback for the next update. GKD develops this principle for
sequence models~\cite{gkd}, while Video-OPD applies it to
student-generated outputs in temporal video grounding~\cite{videoopd}.
Relative-entropy policy improvement further provides a way to turn
scalar evaluations into a distributional target~\cite{reps,mpo}.
Together, these ideas motivate learning from the relative quality of
localizations sampled around the current student prediction instead
of prescribing a single box as its training target.

We study a specialist-guided
STVG-TTA framework built on on-policy spatial distribution
distillation. Around each current student box, we define a local
action distribution and sample candidate localizations. Their
agreement with the fixed Grounding DINO reference supplies a
geometric reward, which reweights the sampled actions into a
feedback-improved distribution. We distill this distribution through
the student's lightweight spatial decoder and regenerate candidate
actions after each update. In this way, the supervision follows the
student's changing spatial predictions. Because the adapted decoder
is shared across frames, observations at a few locations can
influence the predicted trajectory beyond those frames.

The same query-specific feedback presents a different challenge in an
online stream. Two consecutive video--query pairs may involve
different people, objects, and actions even when they share aspects of
the target domain. A correction that benefits one referred instance
may consequently be less useful for the next, while shifts in
appearance or capture conditions may persist across multiple
arrivals. Continual TTA studies have shown that indiscriminate online
updates can accumulate errors or lose previously useful
knowledge~\cite{cotta,rotta}. For STVG, the unit of correction is
particularly instance-dependent: each query identifies a different
spatio-temporal target, but the decoder parameters are shared across
the stream. This motivates distinguishing the model state used to
correct the current instance from the state inherited by subsequent
queries.

Our dual-timescale update implements this distinction with a compact
spatial adaptation interface. A query residual and selected LayerNorm
parameters are optimized together and used in full to predict the
current tube. Afterward, the residual and optimizer state are reset,
while only an attenuated fraction of the LayerNorm displacement is
carried into the initialization for the next arrival. The full local
state therefore serves the current instance, whereas the inherited
state evolves more gradually across the target stream. Both operate
within the original STVG model and its native temporal interval.

We evaluate the framework on bidirectional cross-dataset transfer
between VidSTG and HC-STVG and examine robustness under same-domain
clean inputs and video corruptions at multiple severities. We further
study the roles of specialist feedback, on-policy sampling, decoder
adaptation, and cross-query parameter consolidation, together with
the cost of sparse specialist access.

Our contributions are threefold:
\begin{itemize}
\item We formulate sparse specialist-guided spatial adaptation for STVG,
keeping the temporal interval explicit and distinguishing frame-level
reference agreement from complete language-conditioned tube accuracy.
\item We implement on-policy spatial distribution distillation through a
compact decoder interface and evaluate the fixed system on complete official
cross-domain populations, including complete negative results and matched
direct-regression controls.
\item We provide controlled evidence separating spatial objective, parameter
readout and temporal coverage. Common-state and output-budget controls delimit
OPD's advantage; cached Query/LN diagnostics and temporal--spatial factorial
readouts identify useful mechanisms without claiming universal superiority or
an established benefit from state inheritance.
\end{itemize}

\section{Related Work}
\label{sec:related-work}

\subsection{Spatio-Temporal Video Grounding}

STVG architectures integrate language, object identity, and temporal structure
to predict the trajectory of a referred instance.
The VidSTG work uses spatio-temporal graph reasoning to model object relations
under multi-form language queries~\cite{vidstg}.
STVGBert directly predicts spatial boxes and temporal boundaries through a
visual--linguistic transformer~\cite{stvgb}, while TubeDETR combines multimodal
video--text encoding with space--time decoding~\cite{tubedetr}.
Subsequent architectures strengthen the association of the target across
frames: STCAT introduces a video-level multimodal template~\cite{stcat},
CoSTA couples temporal and spatial predictions through an entangled
query~\cite{costa}, and CG-STVG refines instance context to guide successive
decoding stages~\cite{cgstvg}.
This line of work provides joint tube-prediction architectures on which
test-time spatial adaptation can operate.

Within this joint prediction process, appearance and event dynamics provide
complementary information.
Collaborative static and dynamic streams explicitly model these two sources
and exchange information between them~\cite{dualstream}.
TA-STVG further organizes target-aware reasoning hierarchically, extracting
target-relevant temporal cues before activating fine-grained spatial
attributes~\cite{tastvg}.
This distinction motivates our choice of adaptation interface: the native
temporal prediction defines the event scope, and the existing spatial decoder
receives the specialist-guided update.

Pretrained vision--language models also provide transferable localization
knowledge for STVG.
VideoGrounding-DINO extends image grounding representations to open-vocabulary
video grounding~\cite{videodino}.
E3M performs query-time visual and textual modulation and alternates temporal
relevance estimation with spatial region identification~\cite{e3m}.
The PG-TAF framework introduced with OmniGround pairs temporal grounding by
a multimodal large language model with spatial propagation within the
selected segment~\cite{omniground}.
Our framework uses this access to pretrained spatial knowledge to adapt an
existing STVG model: a frozen specialist supplies reference boxes within the native
interval, geometric overlap evaluates student-generated candidates, and the student's updated decoder
produces the final trajectory.

\subsection{Test-Time Adaptation and Online State}

TTA extracts learning signals from unlabeled target inputs to update a
deployed model.
Tent minimizes prediction entropy through normalization
parameters~\cite{tent}; MEMO minimizes the entropy of the marginal prediction
over augmented views of a test example~\cite{memo}.
Test-time prompt tuning applies confidence selection and augmentation-based
entropy minimization to vision--language prompts~\cite{tpt}.
EATA selects reliable, nonredundant samples and regularizes updates to
preserve source knowledge~\cite{eata}, while SAR combines reliable sample
selection with sharpness-aware entropy minimization~\cite{sar}.
For video models, ViTTA aligns online feature statistics with training
statistics and enforces consistency across temporal
augmentations~\cite{vitta}.
These objectives use confidence, invariance, or distribution alignment as
adaptation signals; their relationship to task accuracy depends on the target
setting~\cite{pitfalls}.
For STVG, we introduce spatial agreement with an external grounding
specialist as feedback about the referred instance's location.

Continual adaptation additionally determines how information from earlier
inputs influences subsequent predictions.
CoTTA combines a weight-averaged teacher with augmentation-averaged
predictions and stochastic restoration of source parameters~\cite{cotta}.
NOTE addresses temporal correlation through instance-aware normalization and
prediction-balanced memory~\cite{note}, while RoTTA combines robust
normalization with a memory bank that considers uncertainty and
timeliness~\cite{rotta}.
ROID averages source and adapted weights to preserve generalization alongside
certainty and diversity weighting and prior correction~\cite{roid}.
HisTPT retains information in local, hard-sample, and global knowledge banks
and retrieves it to regularize test-time prompt updates~\cite{histpt}.
These approaches organize the retention of learned information across a
stream.
We organize the spatial decoder state by its lifetime, using a query-specific
component for the current instance and a gradually consolidated LayerNorm
component to initialize later queries.

\subsection{Language-Conditioned Grounding Specialists}

Open-vocabulary detectors make pretrained localization knowledge accessible
through language.
ViLD transfers image and text representations from a vision--language teacher
to a region-based detector~\cite{vild}.
OWL-ViT combines contrastive image--text pretraining with transformer-based
detection~\cite{owlvit}, and Detic expands detector vocabularies using
image-level supervision~\cite{detic}.
OWLv2 scales this capability through self-training on image--text pairs with
pseudo-box annotations~\cite{owlv2}.
Together, these works establish mechanisms for localizing concepts beyond
the categories with task-specific box annotations.

Grounded pretraining further connects object locations to descriptive
language, which is the interface required by STVG queries.
MDETR conditions detection on raw text through joint image--text
reasoning~\cite{mdetr}, and GLIP unifies object detection and phrase grounding
during pretraining~\cite{glip}.
Grounding DINO combines grounded pretraining with a feature enhancer,
language-guided query selection, and a cross-modality decoder for open-set
localization~\cite{groundingdino}.
Foundation models also support adaptation of the detector itself: Gao
\etal use a prompt-based mean-teacher framework and instance memory for
test-time adaptive object detection~\cite{ttaod}.
In our framework, Grounding DINO remains frozen and returns spatial evidence
at a sparse set of event frames.
Agreement between these observations and student-sampled boxes supplies the
reward used to update the STVG decoder, connecting frame-level specialist
knowledge to the student's full video trajectory.

\subsection{On-Policy Distillation and Compact Updates}

Distillation for structured prediction depends on which model generates
the outputs used for supervision.
Sequence-level knowledge distillation trains a student on sequences produced
by a teacher~\cite{seqkd}.
GKD instead obtains teacher feedback on student-generated sequences,
connecting supervision to the outputs encountered under the student's own
generation process~\cite{gkd}.
We apply this principle to continuous localization actions: candidate boxes
are sampled around the current student prediction and evaluated by
geometric agreement with the fixed specialist reference.
Fresh candidates are drawn after each spatial update, so each distillation
round evaluates the behavior of the updated student.

Relative-entropy policy improvement provides a mechanism for converting
these action evaluations into a target distribution.
REPS constrains information loss during policy improvement~\cite{reps}.
MPO alternates construction of a nonparametric improved distribution with
fitting a parametric policy in an off-policy reinforcement learning
setting~\cite{mpo}; V-MPO develops on-policy policy iteration using a learned
state-value function~\cite{vmpo}.
Our localization objective uses the same principle of reward-based
reweighting, with rewards computed directly from specialist-box IoU.
The resulting target policy is an exponential reweighting of the current box
policy, and forward-KL distillation fits that target through the spatial
decoder state.
This formulation connects sparse specialist observations to a distributional
update of the deployed grounder.

Compact adaptation interfaces specify where this distributional update enters
the pretrained model.
Visual prompt tuning learns input-space prompts while keeping the backbone
fixed~\cite{vpt}, and LoRA introduces trainable low-rank weight updates into
frozen transformers~\cite{lora}.
Within continual TTA, EcoTTA adapts small auxiliary networks and uses
self-distilled regularization against the frozen original
network~\cite{ecotta}.
Our interface consists of a second-pass spatial query residual and selected
LayerNorm affine parameters in the final spatial decoder block, totaling
1,792 trainable parameters.
The complete adapted state generates the current tube; the query residual
and optimizer then reset, while an attenuated LayerNorm displacement is
consolidated into the inherited state.
This separation lets the same compact interface express current-query
specialization and gradual state evolution across the stream.

\section{Methodology}
\label{sec:methodology}

\begin{figure*}[t]
  \centering
  \includegraphics[width=\textwidth]{figures/pipeline.pdf}
  \caption{\textbf{Overview of One World, Two Timescales.}
  TA-STVG predicts a native tube, whose temporal interval determines the frames queried by the frozen Grounding DINO specialist.
  Over $S$ on-policy rounds, student-generated spatial actions are scored by their IoU with specialist boxes, reweighted into $q\propto\pi e^{R/\tau}$, and distilled through $\mathcal L_{\mathrm{OPD}}$ into the query residual $r$ and selected LayerNorm parameters $\lambda$.
  The full adapted state produces the current tube.
  The purple carryover path attenuates the LayerNorm displacement while retaining the inherited state, as specified in \eqref{eq:slow-writeback}; the query residual and optimizer state reset before the next query.}
  \label{fig:pipeline}
\end{figure*}

\subsection{Problem Setup and Overview}
\label{sec:overview}

We consider an online stream of video--query pairs
$x_n=(\widetilde V_n,Q_n)$, where $\widetilde V_n$ is the video observed at
arrival $n$ and $Q_n$ specifies the referred target.
Let $\mathcal T_n$ denote the sampled frames.
A pretrained STVG model predicts a joint tube
$Y_n=(I_n,B_n)$, where $I_n=[s_n,e_n]$ specifies the temporal extent of the
queried event and $B_n=\{b_{n,t}:t\in\mathcal T_n\}$ gives the framewise
trajectory of the referred instance.

Our goal is to adapt this pretrained grounder to an unlabeled target stream
using sparse spatial feedback from a single frozen grounding specialist.
The key difficulty is that such feedback is informative only at a few observed
frames, whereas STVG requires a complete tube over an event. Moreover, the
specialist prediction is itself an external hypothesis and therefore should
guide the student's behavior without replacing the student's grounding
process.

Our key idea is to formulate spatial correction as \emph{on-policy
	distribution distillation}. The STVG model first uses its native temporal
prediction to determine where specialist evidence should be collected
(Sec.~\ref{sec:temporal-scope}). We then turn the student's spatial output into
a local action policy. The student generates localization actions from its
current policy, while IoU with the fixed Grounding DINO box evaluates
these actions and induces a feedback-improved policy. Distilling this improved policy back into a compact
decoder state updates the student's complete spatial trajectory
(Sec.~\ref{sec:spatial-opd}).

The adapted state is finally organized at two timescales
(Sec.~\ref{sec:state-lifecycle}). The complete spatial update is used to
specialize the model to the current query, whereas only a small LayerNorm
component evolves across queries. This design couples strong current-query
adaptation with conservative online state evolution within a single STVG
model.

Figure~\ref{fig:pipeline} summarizes the complete pipeline.


\subsection{Event-Scoped Specialist Feedback}
\label{sec:temporal-scope}

A spatial specialist can only inspect a limited number of video frames.
Querying arbitrary frames from a long untrimmed video would spend this budget
on observations that may lie outside the event described by the language
query. STVG already provides a useful structure for allocating this budget:
its temporal prediction estimates the portion of the video in which the
referred event occurs. We therefore use the student's native temporal
prediction as an \emph{event scope} for acquiring spatial feedback.

For each arrival, the original native tube provides the fixed interval
$I_n^0=[s_n^0,e_n^0)$. The inherited adaptive state provides the
spatial Before prediction; it does not redefine this cached temporal scope.
Spatial updates retain $I_n^0$ in the final output, so an event portion
outside this interval cannot be recovered by this method.
Let $\tau(t)$ denote the physical timestamp of sampled frame
$t\in\mathcal T_n$.
We partition the physical duration of $I_n^0$ into $K$ equal-width bins and
select the available frame nearest to each bin center $c_{n,k}$:
\begin{equation}
	t_{n,k}
	=
	\arg\min_{t\in\mathcal T_n\cap I_n^0}
	\left|\tau(t)-c_{n,k}\right|,
	\qquad k=1,\ldots,K .
	\label{eq:event-sampling}
\end{equation}
The resulting frames provide sparse observations distributed across the
student-predicted event rather than concentrated at a single temporal
location.

For each selected frame $t_{n,k}$, the frozen Grounding DINO specialist takes
the frame and referring expression and returns candidate boxes
$\{(\widetilde b_{n,k,j},c_{n,k,j})\}_{j=1}^{J_{n,k}}$.
After applying a fixed admission rule, we retain the highest-scoring valid
candidate
$\bar b_{n,k}=\widetilde b_{n,k,j^\star}$ with
$j^\star=\arg\max_j c_{n,k,j}$.
We collect these observations as
$\Omega_n=\{(t_{n,k},\bar b_{n,k})\}$.
The fixed admission rule requires mixed entity/phrase confidence at least
0.35 and a margin of at least 0.05 over the next distinct candidate after
duplicate suppression at IoU 0.5. These rules use no target GT.
When no observation is admitted, no spatial optimizer step is taken and
After equals the inherited Before prediction. Before can differ from the
frozen source prediction because preceding arrivals may have changed
LayerNorm state; the query residual and optimizer still reset.

This module converts the native temporal prediction into an efficient
specialist-acquisition region. Its output is not a spatial trajectory; it is a
small set of event-scoped observations that will be used to evaluate the
student's own localization behaviors. The next module therefore addresses the
central spatial problem: how to propagate these sparse observations to the
full STVG tube.


\subsection{On-Policy Spatial Distribution Distillation}
\label{sec:spatial-opd}

Sparse specialist observations supply reference boxes rather than a
complete trajectory. We compare local student-generated actions by their
IoU with each fixed reference, then fit the resulting weighted target.
This interface differs from directly minimizing L1/GIoU to the reference
in its local sampling and iterative reward reweighting. Its value relative
to direct regression is an empirical question evaluated with a matched
Direct control; distributional supervision alone does not establish a
stronger or semantically richer teacher signal.

\paragraph{Lightweight spatial adaptation interface.}
We first expose a compact set of parameters through which spatial feedback can
modify the STVG decoder. For arrival $n$, the adaptive state is
$\phi_n=[r_n,\lambda_n]\in\mathbb R^{1792}$.
Here, $r_n\in\mathbb R^{256}$ is a residual added to the second-pass spatial
query, while $\lambda_n\in\mathbb R^{1536}$ contains the affine parameters of
three LayerNorm operations in the final spatial decoder block.
All other model parameters remain fixed.

The state is initialized with $r_n^{(0)}=0$ and the LayerNorm state inherited
from the previous arrival. Under state $\phi$, the spatial decoder produces
$B_n(\phi)=\{b_{n,t}(\phi):t\in\mathcal T_n\}$.
Because the same adaptive state participates in spatial decoding across the
sampled sequence, updates inferred from a few observed frames can change the
student's trajectory beyond those observations.

\paragraph{Student spatial action policy.}
The next question is how to expose multiple plausible localization behaviors
around the current prediction. We use the deterministic student box as the
center of an explicit local action policy.
For an observed frame $t\in\Omega_n$, let
$b_{n,t}(\phi)\in(0,1)^4$ denote the normalized center--size prediction and
define its unconstrained representation
$\mu_{n,t}(\phi)=\operatorname{logit}(b_{n,t}(\phi))$.
The student policy is
\begin{equation}
	\pi_{\phi,t}(z)
	=
	\mathcal N
	\!\left(
	z;\mu_{n,t}(\phi),\sigma^2 I
	\right).
	\label{eq:spatial-policy}
\end{equation}
An action is mapped back to a valid normalized box through
$a(z)=\operatorname{sigmoid}(z)$.

The logit coordinate system is useful here because it provides an unconstrained
space for symmetric local exploration while preserving valid box coordinates
after the inverse mapping. The deterministic STVG prediction remains the mean
of the policy, so introducing the stochastic policy does not replace the
student's ordinary spatial readout.

At distillation round $s$, we draw $M$ actions around the current policy mean.
Using antithetic Gaussian noise gives
$z_{n,t,m}^{(s)}=\mu_{n,t}^{(s)}+\sigma\epsilon_m$ with
$\epsilon_{m+M/2}=-\epsilon_m$.
Thus each round explores a balanced local neighborhood around the
\emph{current} student prediction.

\paragraph{Specialist-guided policy improvement.}
Having generated student actions, we compare them with the same fixed
specialist observation. For each sampled action, our geometric reward is
$R_{n,t}(z)=\operatorname{IoU}(a(z),\bar b_{n,t})$.
These rewards are used to improve the current student policy under a
KL regularizer:
\begin{equation}
	q_{n,t}^{(s)}
	=
	\arg\max_q
	\left[
	\mathbb E_{z\sim q} R_{n,t}(z)
	-
	\tau
	D_{\mathrm{KL}}
	\!\left(
	q\,\|\,\pi_{\phi_n^{(s)},t}
	\right)
	\right].
	\label{eq:policy-improvement}
\end{equation}
This objective balances two goals: moving probability mass toward actions
favored by the specialist while retaining the behavioral support of the
current student. It admits the closed-form solution
\begin{equation}
	q_{n,t}^{(s)}(z)
	=
	\frac{
		\pi_{\phi_n^{(s)},t}(z)
		\exp\!\left(R_{n,t}(z)/\tau\right)
	}{
		Z_{n,t}^{(s)}
	} .
	\label{eq:improved-policy}
\end{equation}

Since the candidate actions are sampled from
$\pi_{\phi_n^{(s)},t}$ itself, their normalized feedback weights are
$w_{n,t,m}^{(s)}
\propto
\exp(R_{n,t}(z_{n,t,m}^{(s)})/\tau)$.
The resulting $q_{n,t}^{(s)}$ can therefore be interpreted as a
\emph{feedback-improved version of the student's own spatial policy}: the
student determines which localization behaviors are explored, and the
specialist determines how probability should be redistributed among them.


\paragraph{On-policy distribution distillation.}
The feedback-improved policy $q_{n,t}^{(s)}$ encodes the
specialist's preference over localization actions generated
by the current student. To transfer this feedback into the
spatial decoder, we project the improved policy onto the
student's parameterized Gaussian policy family by minimizing
the forward KL divergence:
\begin{equation}
	\mathcal L_{\mathrm{OPD}}^{(s)}(\phi)
	=
	\frac{1}{|\Omega_n|}
	\sum_{t\in\Omega_n}
	D_{\mathrm{KL}}
	\left(
	\operatorname{sg}[q_{n,t}^{(s)}]
	\,\|\, \pi_{\phi,t}
	\right).
	\label{eq:opd-objective}
\end{equation}

We optimize this objective using the actions sampled from
the current student policy and their specialist-derived
importance weights. With the improved policy held fixed,
the forward-KL gradient is estimated through the weighted
student log-likelihood. Our implementation uses an
analytically centered likelihood estimator, which yields
the same parameter gradient as the sample-weighted
likelihood under the fixed-covariance Gaussian policy
and antithetic sampling. The resulting gradients update
the spatial query residual and decoder LayerNorm parameters,
transferring the specialist's relative preferences into
the student's localization behavior.

For fixed action samples and detached weights, the sample-weighted
Gaussian likelihood has the equivalent weighted-regression form
\begin{equation}
\mathcal L_{\mathrm{sample}}^{(s)}(\phi)
=\frac{1}{2\sigma^2|\Omega_n|}
\sum_{t\in\Omega_n}\left\|\mu_{n,t}(\phi)
-\sum_m w_{n,t,m}^{(s)}\operatorname{sg}[z_{n,t,m}^{(s)}]\right\|_2^2+C.
\label{eq:weighted-regression}
\end{equation}
Here $C$ is independent of the updated mean. The distinction from Direct
therefore lies in the locally constructed target and its resampling across
rounds, rather than escaping regression mathematically. If all sampled IoU
rewards are zero at one observation, its weights are uniform and its
antithetic centered correction vanishes at the sampling mean. This does
not imply a query-wide no-op: other observations and shared decoder state
can still move its box. Zero reward can limit both correction toward an
accurate distant reference and movement toward a poor reference.

The distillation target evolves with the student.
After each optimization step, the updated spatial decoder
induces a new action policy, from which fresh localization
actions are sampled and evaluated against the specialist
evidence. This alternating process of policy improvement
and distillation progressively adapts the student's spatial
predictions within its evolving action distribution. After
$S$ rounds, we use the final adapted state
$\phi_n^\star=\phi_n^{(S)}$ to generate the spatial trajectory
over the native event interval.



\subsection{Query Reset and LayerNorm Writeback}
\label{sec:state-lifecycle}

The spatial state optimized above is tailored to the current referred instance,
but online adaptation also determines the initialization from which future
queries will be processed. The implementation assigns these roles different lifetimes: both parameter
groups form the current output, while the query residual and optimizer reset
before the next arrival. A small LayerNorm displacement is retained by the
fixed writeback rule. Its utility is an empirical question for matched
state-lifecycle controls.

We therefore decompose the adapted state
$\phi_n^\star=[r_n^\star,\lambda_n^\star]$ by lifetime.
For the current query, both components are used without attenuation.
Let $B_n^\star=B_n(\phi_n^\star)$; the final STVG output is
$Y_n^{\mathrm{out}}=(I_n^0,B_n^\star|_{I_n^0})$.
This fast state allows the current prediction to exploit the complete spatial
adaptation capacity.

After producing the current tube, the query residual is reset.
Only the LayerNorm component contributes to the state of the next arrival.
If $\lambda_n^G$ denotes the inherited LayerNorm state before adaptation, the
slow state evolves as
\begin{equation}
	\lambda_{n+1}^{G}
	=
	(1-\alpha)\lambda_n^{G}
	+
	\alpha\lambda_n^\star .
	\label{eq:slow-writeback}
\end{equation}
The optimizer is reinitialized together with the query residual at the next
arrival.

This lifecycle gives the two parameter groups distinct functions.
The query residual supplies 256 local coordinates and the selected
LayerNorm parameters supply 1,536 shared coordinates. This describes the
implemented state lifecycle; it does not establish that both groups are
necessary or that retaining LayerNorm state causally improves grounding.

% Standalone Experiments section. Requires amsmath, graphicx, and booktabs,
% which are already loaded by the user's existing CVPR main.tex.
% The following names are safe in text, tables, captions, and math mode.
\providecommand{\mvIoU}{\ensuremath{\mathrm{m}_{\mathrm{vIoU}}}}
\providecommand{\mtIoU}{\ensuremath{\mathrm{m}_{\mathrm{tIoU}}}}
\providecommand{\msIoU}{\ensuremath{\mathrm{m}_{\mathrm{sIoU}}}}

\section{Experiments}
\label{sec:experiments}

We evaluate our method under natural cross-domain shifts and
same-domain video corruptions. Our experiments investigate
whether sparse specialist feedback improves cross-domain grounding,
how on-policy distillation converts this feedback into spatial
corrections, and how adaptation to the current query interacts
with information retained across an online stream.

\subsection{Experimental Setup}
\label{sec:experimental-setup}

\paragraph{Datasets and protocols.}
We conduct experiments on VidSTG~\cite{vidstg} and HC-STVG-v2,
which associate language queries with temporal event intervals
and spatial object trajectories. TA-STVG~\cite{tastvg} serves as
the pretrained STVG backbone. Our primary setting is bidirectional
cross-domain transfer: VidSTG $\rightarrow$ HC-STVG-v2 and
HC-STVG-v2 $\rightarrow$ VidSTG. We evaluate all official target
queries, comprising 3,482 queries from 237 parent videos on
HC-STVG-v2 and 10,303 queries from 732 videos on VidSTG.
Predictions are obtained under three fixed stream orders, with
the adaptive state initialized independently for each order.
No target-domain annotations are used for adaptation.

For same-domain robustness, we use the corresponding in-domain
checkpoint and one fixed query per parent video (237 queries
on HC-STVG-v2 and 732 on VidSTG). We evaluate clean videos
and five types of transient corruption: frame dropping, frame
freezing, motion blur, occlusion, and exposure changes. Each
corruption affects 2.5\%, 5\%, or 10\% of the video's physical
duration. Each condition begins from an independently initialized
source model.

\paragraph{Evaluation metrics.}
Following standard STVG evaluation, we report mean video IoU
(\mvIoU), temporal IoU (\mtIoU), and spatial IoU (\msIoU),
along with the percentage of queries whose video IoU exceeds
0.3 and 0.5 (R@0.3 and R@0.5). The full cross-domain table
uses query-level averages over the three stream orders. For
paired comparisons, we first average queries and orders within
each parent video and then use 10,000 paired parent-level
bootstrap resamples to obtain 95\% confidence intervals.

\paragraph{Implementation details.}
We use a single frozen Grounding DINO~\cite{groundingdino}
to provide frame-level references. Unless specified otherwise,
we uniformly select $K=4$ frames within the native predicted
event interval and use the highest-scoring admitted reference
at each selected frame. The student samples $M=32$ antithetic
Gaussian localization actions per observed frame and round.
Only a 256-dimensional spatial query residual and 1,536
LayerNorm affine parameters are trainable, totaling 1,792
parameters. For the VidSTG target, we use learning rate
$0.03$, exploration scale $\sigma=0.1$, temperature $\tau=0.25$,
10 distillation rounds, and LayerNorm consolidation factor
$\alpha=1/8$. For the HC-STVG-v2 target, the corresponding
values are $0.01$, $0.025$, $0.05$, 40, and $1/16$.
The settings are fixed using development subsets of 32 parent
videos per target; no evaluation labels are used during an
adaptation stream. We separately report analyses excluding these
parents. The remaining sources have historical evaluation exposure
and are not described as a newly untouched confirmation population. The final round determines the current
prediction, while the native temporal interval remains unchanged.

\subsection{Cross-Domain Grounding}
\label{sec:cross-domain-results}

\paragraph{Motivating spatial failures.}
We first examine source-trained TA-STVG and
TubeDETR~\cite{tubedetr} predictions under the two cross-domain
transfers. Even when the predicted temporal interval has
$tIoU>0.5$, spatial localization can remain incorrect.
Among such temporally successful predictions, 38.0\% of
TA-STVG and 43.6\% of TubeDETR outputs on HC-STVG-v2 have
spatial IoU at most 0.5. The corresponding rates on VidSTG
are 55.3\% and 34.8\%. These conditional results identify a
persistent spatial correction opportunity after an event has
been located, motivating the use of the native interval as
the scope for acquiring additional spatial evidence.

\paragraph{Main results.}
Table~\ref{tab:cross-domain-main} reports the full official
cross-domain evaluation. We compare with Source Only and
STVG adaptations of Tent~\cite{tent} and SAR~\cite{sar},
and parameter-free DINO Refine using exactly the same four
admitted references. DINO Refine interpolates absolute center--size
boxes between admitted anchors with nearest endpoint extension;
it is not an update at only four frame positions.
An in-domain supervised model is shown separately as a
reference trained with target-domain annotations.

\begin{table*}[t]
	\centering
	\caption{Cross-domain STVG results on all official target queries.
		All values are percentages. Tent/SAR here retain the original
        evaluated configurations; separately calibrated configurations are
        not substituted before their complete formal evaluation.
        The in-domain supervised reference
		uses target-domain training labels and is not a TTA baseline.
		Bold indicates the best result among methods without target
		training annotations; ties are left unbolded.}
	\label{tab:cross-domain-main}
	\setlength{\tabcolsep}{8pt}
	\small
	\begin{tabular}{lccccc}
		\toprule
		Method & \mvIoU{} $\uparrow$ & \mtIoU{} $\uparrow$ &
		\msIoU{} $\uparrow$ & R@0.3 $\uparrow$ & R@0.5 $\uparrow$ \\
		\midrule
		\multicolumn{6}{l}{\textit{VidSTG $\rightarrow$ HC-STVG-v2}} \\
		\midrule
		Source Only         & 20.017 & 40.311 & 48.550 & 27.570 & 7.639 \\
		Tent-STVG           & 16.598 & 33.333 & 48.536 & 20.132 & 5.035 \\
		SAR-STVG            & 19.939 & 40.152 & 48.556 & 27.465 & 7.486 \\
		DINO Refine          & 20.673 & 40.311 & 43.631 & 28.920 & 9.075 \\
		\textbf{Ours}       & \textbf{22.583} & 40.311 & \textbf{53.625} &
		\textbf{32.970} & \textbf{11.784} \\
		\midrule
		In-domain (supervised) & 29.056 & 57.280 & 49.332 & 45.951 & 14.302 \\
		\midrule
		\multicolumn{6}{l}{\textit{HC-STVG-v2 $\rightarrow$ VidSTG}} \\
		\midrule
		Source Only         & 11.129 & 31.703 & 32.758 & 12.831 & 3.300 \\
		Tent-STVG           & 8.721 & 24.973 & 32.767 & 9.774 & 2.271 \\
		SAR-STVG            & 11.129 & 31.703 & 32.758 & 12.831 & 3.300 \\
		DINO Refine          & 13.985 & 31.703 & 35.710 & \textbf{18.024} & 6.920 \\
		\textbf{Ours}       & \textbf{14.097} & 31.703 & \textbf{39.855} &
		17.904 & \textbf{7.037} \\
		\midrule
		In-domain (supervised) & 19.296 & 42.750 & 42.325 & 28.137 & 13.103 \\
		\bottomrule
	\end{tabular}
\end{table*}

Our method improves whole-tube localization in both directions.
On HC-STVG-v2, \mvIoU{} increases from 20.017\% to 22.583\%,
with gains of 5.400 and 4.145 percentage points at R@0.3
and R@0.5. On VidSTG, \mvIoU{} increases from 11.129\%
to 14.097\%, while R@0.5 increases from 3.300\% to 7.037\%.
The gains arise from spatial correction: \mtIoU{} is unchanged,
while \msIoU{} increases from 48.550\% to 53.625\% on
HC-STVG-v2 and from 32.758\% to 39.855\% on VidSTG.
Thus, a small number of frame-level references can improve
localization across the predicted event through parameter
updates to the shared spatial decoder.

Excluding the 32 development parent videos per target,
the paired parent-level improvement over Source Only is
$+2.409$ percentage points (95\% CI: $[1.891,2.907]$)
on the remaining 205 HC-STVG-v2 parents and $+3.329$
points ($[3.017,3.646]$) on the remaining 700 VidSTG
parents. The improvements remain positive after this
exclusion, supporting the generality of the gains within
the evaluated cross-domain populations.

\subsection{Robustness to Video Corruptions}
\label{sec:robustness}

Natural domain shifts change the distribution of scenes
and queries, whereas transient corruption disrupts the
video evidence within the same domain. We examine whether
spatial adaptation remains useful under these localized
changes. Table~\ref{tab:robustness} reports the clean
results, individual corruption families at 5\% coverage,
and means across the five families at each severity level.

\begin{table*}[t]
	\centering
	\caption{Same-domain robustness under transient video
		corruptions. The first block reports clean performance and
		individual corruption types at 5\% physical-time coverage.
		The second block averages the five types at each severity.
		All entries are \mvIoU{} (\%); $\Delta$ is measured in
		percentage points.}
	\label{tab:robustness}
	\setlength{\tabcolsep}{7pt}
	\small
	\begin{tabular}{lcccccc}
		\toprule
		& \multicolumn{3}{c}{HC-STVG-v2} & \multicolumn{3}{c}{VidSTG} \\
		\cmidrule(lr){2-4}\cmidrule(lr){5-7}
		Condition & Frozen & Ours & $\Delta$ & Frozen & Ours & $\Delta$ \\
		\midrule
		Clean            & 30.684 & 34.450 & +3.766 & 21.224 & 23.679 & +2.456 \\
		Frame Drop       & 29.961 & 33.490 & +3.529 & 19.257 & 21.295 & +2.039 \\
		Frame Freeze     & 30.631 & 34.555 & +3.925 & 21.269 & 23.357 & +2.088 \\
		Motion Blur      & 30.524 & 33.617 & +3.093 & 20.979 & 23.009 & +2.030 \\
		Occlusion        & 30.400 & 33.432 & +3.033 & 20.875 & 23.077 & +2.201 \\
		Exposure         & 30.536 & 34.347 & +3.811 & 21.208 & 23.345 & +2.137 \\
		\midrule
		\multicolumn{7}{l}{\textit{Mean across corruption families}} \\
		\midrule
		2.5\% coverage   & 30.698 & 34.230 & +3.532 & 20.931 & 23.215 & +2.284 \\
		5\% coverage     & 30.410 & 33.888 & +3.478 & 20.718 & 22.817 & +2.099 \\
		10\% coverage    & 29.807 & 33.347 & +3.540 & 20.183 & 22.361 & +2.178 \\
		\bottomrule
	\end{tabular}
\end{table*}

Our method improves \mvIoU{} on clean inputs and under
every evaluated corruption. At 5\% coverage, improvements
range from 3.033 to 3.925 percentage points on HC-STVG-v2
and from 2.030 to 2.201 on VidSTG. Across the full set of
sixteen clean and corrupted conditions per target, all
paired parent-bootstrap intervals for Ours minus Frozen
are positive.

The gains also persist as the affected duration increases.
Averaged over corruption families, the improvements at
10\% coverage are 3.540 and 2.178 percentage points on
HC-STVG-v2 and VidSTG, respectively. This complements
the cross-domain evaluation by showing that the spatial
updates remain useful under changes to the incoming video
signal without updating the pretrained encoders.

\subsection{Understanding On-Policy Spatial Distillation}
\label{sec:opd-ablation}

Our spatial objective uses rewards assigned to actions
sampled from the current student policy, then refreshes
those actions after each update. We isolate the role of
these choices with three controlled variants. \emph{Direct}
minimizes L1/GIoU (weights 5/2) to the same admitted references,
using the same trainable interface, rounds and stream lifecycle. \emph{Shuffled
	Feedback} preserves the sampled actions and reward values
but randomizes their correspondence. \emph{Fixed Rollout}
retains candidates drawn from the initial student policy
rather than refreshing them during adaptation. Both variants
use the same specialist observations and trainable interface.

\begin{table*}[t]
	\centering
	\caption{Effect of feedback correspondence and on-policy
		sampling. Entries are paired \mvIoU{} differences (percentage
		points) between the full method and each variant, with 95\%
		parent-bootstrap confidence intervals. Cross-domain rows
		use clean inputs; same-domain rows average five corruption
		families at 5\% coverage on fixed 128-parent cohorts.}
	\label{tab:opd-ablation}
	\setlength{\tabcolsep}{5pt}
	\small
	\begin{tabular}{lccc}
		\toprule
		Setting & Full $-$ Direct & Full $-$ Shuffled & Full $-$ Fixed Rollout \\
		\midrule
		VidSTG $\rightarrow$ HC-STVG-v2
		& $+0.576\;[-0.236,1.433]$ & $+3.837\;[2.257,5.335]$ & $+1.124\;[-0.233,2.244]$ \\
		HC-STVG-v2 $\rightarrow$ VidSTG
		& $-0.213\;[-0.625,0.216]$ & $+8.820\;[6.861,10.979]$ & $+3.860\;[2.415,5.334]$ \\
		HC-STVG-v2 (corrupt)
		& $+0.603\;[-1.069,2.267]$ & $+6.320\;[4.595,7.991]$ & $+3.357\;[1.942,4.691]$ \\
		VidSTG (corrupt)
		& $+0.523\;[-0.343,1.494]$ & $+10.222\;[7.455,13.033]$ & $+2.579\;[0.579,4.487]$ \\
		\bottomrule
	\end{tabular}
\end{table*}

\paragraph{Direct regression.}
The Full-minus-Direct intervals in Table~\ref{tab:opd-ablation}
span zero in all four settings. The paired means are $+0.576$
points on cross-domain HC-STVG-v2 and $-0.213$ on cross-domain
VidSTG. These comparisons do not establish stable average OPD
superiority; reward correspondence and policy-refresh ablations test
specific choices within OPD rather than replace this direct comparator.

\paragraph{Feedback correspondence.}
Table~\ref{tab:opd-ablation} shows that the full method
outperforms Shuffled Feedback in all four settings, with
paired gains of 3.837--10.222 percentage points. Since
shuffling retains the reward distribution but destroys its
association with the sampled localization actions, these
results establish that useful feedback depends on assigning
higher weight to the appropriate student-generated actions.

\paragraph{Policy refresh.}
Refreshing the action distribution also improves grounding
in several settings. Compared with Fixed Rollout, the full
method gains 3.860 points on cross-domain VidSTG and
3.357 and 2.579 points on the same-domain corruption
panels. This shows that reevaluating actions around the
updated prediction can make feedback more effective than
reusing the initial action distribution. The corresponding
cross-domain HC-STVG-v2 contrast has a positive mean but
a confidence interval overlapping zero.

\subsection{Fast Adaptation and Slow State Evolution}
\label{sec:state-ablation}

We next examine which parameters carry spatial correction
and how their updates influence subsequent queries.
\emph{Query-only} adapts the 256-dimensional query residual;
\emph{LN-only} adapts the 1,536 LayerNorm affine parameters.
The joint variant with $\alpha=0$ adapts both groups for
the current query but performs no cross-query writeback.
The full method uses both groups locally and retains a
fraction of the LayerNorm displacement across queries.

\begin{table*}[t]
	\centering
	\caption{Parameter and state ablations. Entries are paired
		\mvIoU{} differences (percentage points) between the full
		method and each variant, with 95\% parent-bootstrap
		confidence intervals. The cohorts and settings match
		Table~\ref{tab:opd-ablation}.}
	\label{tab:state-ablation}
	\setlength{\tabcolsep}{6pt}
	\small
	\begin{tabular}{lccc}
		\toprule
		Setting & Full $-$ Query-only & Full $-$ LN-only & Full $-$ $\alpha=0$ \\
		\midrule
		VidSTG $\rightarrow$ HC-STVG-v2
		& $+1.182\;[-0.070,2.247]$ & $-0.170\;[-0.734,0.296]$ & $+0.238\;[0.011,0.488]$ \\
		HC-STVG-v2 $\rightarrow$ VidSTG
		& $+3.322\;[2.048,4.608]$ & $+0.464\;[0.090,0.909]$ & $+0.258\;[-0.182,0.691]$ \\
		HC-STVG-v2 (corrupt)
		& $+3.942\;[2.686,5.152]$ & $+0.009\;[-0.490,0.468]$ & $+0.217\;[-0.279,0.673]$ \\
		VidSTG (corrupt)
		& $+1.438\;[-0.676,3.415]$ & $+0.420\;[-0.034,0.936]$ & $+0.480\;[-0.257,1.211]$ \\
		\bottomrule
	\end{tabular}
\end{table*}

As shown in Table~\ref{tab:state-ablation}, the two
parameter groups contribute differently across datasets.
On cross-domain VidSTG, the full method exceeds Query-only
by 3.322 points and LN-only by 0.464 points. On
HC-STVG-v2, LN-only captures much of the full gain,
particularly under same-domain corruption. The contrast
indicates that the LayerNorm parameters are an important
adaptation interface, while the query residual has a positive contrast in this particular
VidSTG panel. Its indispensable contribution is not established: the enlarged
matched state controls and independent LN-only stream remain outstanding.

We also analyze the state inherited from preceding queries.
Let $v_F$, $v_B$, and $v_A$ denote performance from the
frozen source state, the inherited state before adapting
the current query, and the fully adapted current state.
Then $v_A-v_F=(v_B-v_F)+(v_A-v_B)$. Excluding development
sources from the full cross-domain population, the inherited
component $v_B-v_F$ is $+2.054$ points on HC-STVG-v2
and $+1.699$ on VidSTG, while the current-query component
$v_A-v_B$ is $+0.355$ and $+1.629$, respectively. Thus,
performance depends on both the current correction and
the state from which it begins, with substantially different
balances across targets.

The $\alpha=0$ intervention directly tests the effect of
cross-query consolidation. Its difference from the full
method is positive on cross-domain HC-STVG-v2
($+0.238$ points, 95\% CI $[0.011,0.488]$); the other
three comparisons have intervals spanning zero. The
inherited-state decomposition and this intervention answer
different questions: the former describes the state of the
full online model, whereas the latter isolates the effect
of the writeback rule.

\subsection{Specialist Observation Budget}
\label{sec:observation-budget}

We vary the number of uniformly sampled event frames,
$K\in\{1,2,4,8\}$, while keeping the remaining configuration
fixed. Table~\ref{tab:budget} summarizes the resulting
cross-domain improvements and recorded decoder-fitting
costs on the fixed 128-parent cohorts.

\begin{table*}[t]
	\centering
	\caption{Specialist observation budget on cross-domain
		clean inputs. Improvements are parent-averaged \mvIoU{}
		percentage points over Frozen, with 95\% paired confidence
		intervals. Fit time is measured per arrival and excludes
		shared capture and specialist inference.}
	\label{tab:budget}
	\setlength{\tabcolsep}{9pt}
	\small
	\begin{tabular}{ccccc}
		\toprule
		& \multicolumn{2}{c}{HC-STVG-v2 target} & \multicolumn{2}{c}{VidSTG target} \\
		\cmidrule(lr){2-3}\cmidrule(lr){4-5}
		$K$ & $\Delta\mvIoU$ (pp) & Fit (s) & $\Delta\mvIoU$ (pp) & Fit (s) \\
		\midrule
		1 & $+0.655\;[-0.791,1.928]$ & 1.326
		& $+2.249\;[0.985,3.483]$ & 0.322 \\
		2 & $+0.664\;[-0.739,1.859]$ & 1.351
		& $+3.649\;[2.205,5.081]$ & 0.334 \\
		4 & $+1.346\;[-0.046,2.526]$ & 1.356
		& $+3.871\;[2.430,5.290]$ & 0.347 \\
		8 & $+0.342\;[-1.393,1.884]$ & 1.413
		& $+4.458\;[2.847,6.078]$ & 0.343 \\
		\bottomrule
	\end{tabular}
\end{table*}

The method obtains improvements from sparse observations.
On VidSTG, all four budgets have positive paired confidence
intervals relative to Frozen; increasing $K$ from one to
four increases the mean gain from 2.249 to 3.871 points.
The differences between $K=2$ and $K=4$, and between
$K=4$ and $K=8$, are less conclusive. On HC-STVG-v2,
$K=4$ yields the highest mean gain among the tested
budgets, while the additional observations at $K=8$
do not improve the mean. The main four-frame configuration was fixed before this
budget evaluation; these results do not select a new test-set budget.

A specialist reference is acquired once per selected
frame and reused across distillation rounds. Consequently,
the optimization cost is separate from the cost of
specialist observations. For the default $K=4$, the
recorded fitting stage takes 1.356 seconds per arrival
on HC-STVG-v2 and 0.347 seconds on VidSTG. These are
fit-stage measurements, not full end-to-end inference
times.

\subsection{Analysis of Spatial Corrections}
\label{sec:correction-analysis}

\paragraph{Corrections beyond observed frames.}
We separate admitted observed frames, unobserved real decoder frames
inside Native WHEN, interpolation-only frames inside Native, and
annotated event frames outside Native. On the matched P2 population
excluding development parents, the current After-minus-Before change
on unobserved real decoder frames is $+5.419$ points on cross-domain
VidSTG (95\% parent CI $[2.594,8.278]$) and $-0.581$ on
cross-domain HC-STVG-v2 ($[-3.238,1.865]$). The corresponding
After-minus-DINO differences are $+2.647$ ($[-0.146,5.777]$)
and $+4.394$ ($[0.193,8.450]$). These contrasts distinguish the
current spatial update from the complete online model's advantage
over parameter-free refinement. Common observed/unobserved parent
intersections are reported separately: for HC, the latter interval
then spans zero. Interpolation alone is not evidence of a learned
decoder change, and geometric changes outside Native WHEN do not
recover temporal coverage.

\paragraph{Expert quality and local actionability.}
Using the actual saved rewards, we stratify expert frames by
expert--GT IoU $Q\geq0.5$ or $Q<0.5$ and by whether any first-round
action has positive reward. In the low-$Q$/all-zero strata,
Full-minus-Direct current frame-IoU changes are $+11.520$ points
on cross-domain HC ($[4.066,19.682]$, 22 parents) and $+9.751$
under HC corruption ($[1.474,18.744]$, 32 parents). The analogous
Vid intervals span zero. However, high-$Q$/all-zero frames favor
Direct in the HC comparisons, by $29.329$ and $40.875$ points on
only 5 and 11 parents. These descriptive geometric strata suggest
a conditional conservative-update tradeoff, not a universal
reliability advantage or verified identity-error detector.
Full and Direct have their own inherited Before states.
All-zero first-round reward occurs in 1.274--6.593\% of the
high-$Q$ observations across the four settings. It remains a
correction failure mode, even though the predeclared resource
trigger for a changed-reward pilot is not met.

\paragraph{Qualitative behavior.}
Aligned Native/Before/After/GT examples show both improvements and
regressions, with the four real specialist observations and sampled
intermediate frames marked. Specialist agreement can increase while
GT overlap decreases. Geometric disagreement alone does not establish
an identity switch; such an attribution requires instance evidence
and inspection of the actual sequence.

\paragraph{Offline spatial and temporal headroom.}
On the matched 128-parent cross-domain populations, the Native,
GT-space/Native-time, and Native-space/GT-time video IoUs are
22.156, 42.239, and 51.417\% for HC, and 14.466, 36.614,
and 38.213\% for VidSTG. The temporal oracle uses only the
original native trajectory and its supported interpolation;
unpredicted frames are never filled with GT spatial boxes.
Temporal headroom is at least as large as spatial headroom in
these settings, so the motivating spatial errors do not imply a
universal dominant spatial bottleneck. Native WHEN covers about
56\% of annotated event support. These offline interventions are
diagnostic ceilings, not deployable comparisons or an additive
decomposition of task error.

\paragraph{Complementary temporal corrections.}
We finally inspect temporal refinement with spatial boxes
held fixed, using a separate diagnostic set of 32 parent
videos per target. An existing actionness-based temporal
projection improves \mvIoU{} over the native interval by
$+3.755$ points on HC-STVG-v2 and $+1.634$ on VidSTG.
An offline ground-truth-guided temporal-head diagnostic
reveals further capacity for correction, particularly
on VidSTG. These Native-space controls demonstrate deployable temporal headroom.
Their interaction with spatial adaptation requires a matched factorial comparison.
Our main method concentrates on spatial correction
within the native event scope. Joint temporal adaptation is not
part of the evaluated primary method.


	
	
\subsection{Common-State Spatial and Parameter-Mean Controls}
The new common-state study uses every original matched cell, preserving Full's
initial state, observations and saved reference. All four primary Full-minus-
Direct intervals span zero. Full-minus-Capped-Direct is negative in both Vid
settings: cross-domain $-0.410\;[-0.718,-0.117]$ pp and same-domain corruption
$-0.303\;[-0.620,-0.006]$ pp. The cap bounds output movement; it does not force
exact norm equality. These controls are conditional on Full history rather
than three independent online streams.

The corrected Native-support Query observability study uses all 64 existing
development parents. Query-observable parameter-mean gains are $+0.002$ pp on
HC and $-0.004$ pp on Vid, while HC Query-minus-LN mean is
$-0.147\;[-0.238,-0.072]$ pp. Candidate oracle and expert-top1 readouts are
distinct from the deployed true parameter mean. Complete signed-pair replay
finds median reward-margin-to-temperature ratios of 0.02056 and 0.01162,
corresponding to 1.028\% and 0.581\% pair retention. Unequal positive/negative
amplitudes contribute a midpoint term; the equal-amplitude tanh illustration
alone is insufficient. These observations support mean-readout shrinkage in
this scope without identifying its unique causal role or selecting a new
temperature. Full Query-level OPD is not introduced.

\subsection{Matched Temporal--Spatial Readouts}
We cross frozen Native WHEN (T0) and the unchanged existing Actionness
Projection (T1) with frozen Native, Full OPD, common-state Direct and Capped
Direct. Exact cache matching admits all original P6 parents (32 per target)
and both existing orders, totaling 128 Common-State cells. The primary panel
uses original order1; the secondary clusters both orders within the same
parents. The other 1,664 cells lack complete corresponding temporal captures
and are not represented as a full extension. Every readout is sealed before
offline scoring; missing decoder support has zero IoU and no GT completion.

\begin{table}[t]
\centering\small
\caption{Primary cached time--space contrasts (pp), 32 historical parents per target; 10,000 paired parent bootstrap, unadjusted 95\% CI.}\label{tab:time-space}
\begin{tabular}{lcc}\toprule
Contrast & HC & Vid \\
\midrule
OPD T1 $-$ T0 & $+3.823\;[+0.336,+7.169]$ & $+2.079\;[+0.269,+4.259]$ \\
OPD/Native interaction & $+0.068\;[-0.794,+0.905]$ & $+0.445\;[-0.663,+1.479]$ \\
OPD $-$ cap at T1 & $+0.397\;[-0.574,+1.404]$ & $-0.472\;[-1.148,+0.201]$ \\
\bottomrule\end{tabular}\end{table}
Actionness Projection improves the OPD readout in both targets, but neither
primary interaction interval excludes zero. OPD-minus-cap at projected time
also spans zero in both primary panels; the two-order Vid comparison is
$-0.771\;[-1.680,-0.020]$ pp. Mean annotated event coverage increases by
19.146 and 5.315 pp, with 1,087 and 1,275 newly included annotated frames
outside Native and no added frames outside decoder support. Four HC and one
Vid parents lose more than 5 pp under OPD's temporal change. Fixed-event
spatial IoU stays unchanged for each spatial readout. This supports useful
temporal coverage recovery without establishing an OPD-specific joint
mechanism or an independently validated new temporal adaptation method.

\section{Conclusion}
Sparse specialist feedback can improve a compact STVG spatial adaptation
system under the fixed cross-domain protocol. The current evidence supports
the complete system gain and distinguishes temporal coverage, spatial
objective and parameter readout; stable OPD-over-Direct superiority,
indispensable Query residuals and causal gains from inherited state remain
unestablished. Enlarged state, calibrated baseline and full deployment timing
claims require completion of their original locked evaluations. We retain the
fixed Full method and report conditional negative controls rather than
promote a new Query or temporal method from these diagnostics.


	{\small
		\bibliographystyle{ieeenat_fullname}
		\bibliography{references}
	}
\end{document}
